\chapter*{PHỤ LỤC: BẢNG THUẬT NGỮ}
\addcontentsline{toc}{chapter}{Phụ lục: Bảng thuật ngữ}

Dưới đây là các thuật ngữ kỹ thuật thường gặp trong lĩnh vực YOLO và Object Detection, sắp xếp theo nhóm chủ đề.

\section*{Kiến trúc mạng (Architecture)}

\begin{longtable}{p{3.5cm}p{12cm}}
\toprule
\textbf{Thuật ngữ} & \textbf{Giải thích} \\
\midrule
\endfirsthead
\toprule
\textbf{Thuật ngữ} & \textbf{Giải thích} \\
\midrule
\endhead

\textbf{Backbone} & Phần ``xương sống'' của model, trích xuất features từ ảnh đầu vào. Ví dụ: Darknet, CSPDarknet, ResNet, EfficientNet. \\
\midrule
\textbf{Neck} & Kết nối Backbone và Head, kết hợp features ở nhiều scales. Ví dụ: FPN, PAN, BiFPN. \\
\midrule
\textbf{Head} & Phần cuối cùng, dự đoán bounding boxes và class probabilities. Có 2 loại: coupled (chung) và decoupled (tách riêng cls/reg). \\
\midrule
\textbf{CNN} & Convolutional Neural Network --- mạng neural dùng phép tích chập (convolution) để trích xuất features từ ảnh. Mỗi layer học nhận diện patterns (cạnh, góc, hình dạng). \\
\midrule
\textbf{Convolution} & Phép tính trượt 1 bộ lọc (kernel) qua ảnh, tạo feature map. Kernel $3\times3$ nhìn vùng $3\times3$ pixels. \\
\midrule
\textbf{Depthwise Conv} & Convolution áp dụng riêng cho từng channel (thay vì tất cả channels cùng lúc) --- giảm computation đáng kể. \\
\midrule
\textbf{Residual Connection} & Đường tắt (shortcut) nối input trực tiếp đến output: $y = F(x) + x$. Giải quyết vanishing gradient khi mạng rất sâu (ResNet). \\
\midrule
\textbf{CSP} & Cross Stage Partial --- chia features thành 2 phần: 1 phần qua processing, 1 phần bypass $\rightarrow$ giảm 20\% computation. \\
\midrule
\textbf{FPN} & Feature Pyramid Network --- kết hợp features từ nhiều scales (top-down pathway) để detect vật ở mọi kích thước. \\
\midrule
\textbf{PAN / PANet} & Path Aggregation Network --- thêm bottom-up pathway vào FPN $\rightarrow$ localization signals mạnh hơn. \\
\midrule
\textbf{SPP / SPPF} & Spatial Pyramid Pooling --- MaxPool nhiều kích thước $\rightarrow$ tăng receptive field. SPPF: phiên bản nhanh hơn (nối tiếp thay song song). \\
\midrule
\textbf{Re-parameterization} & Kỹ thuật training multi-branch, inference single-branch. Gộp nhiều Conv thành 1 Conv duy nhất bằng phép cộng ma trận. \\
\midrule
\textbf{Attention} & Cơ chế cho model ``tập trung'' vào vùng quan trọng. Self-attention so sánh mọi vị trí với nhau (global receptive field). \\
\midrule
\textbf{FlashAttention} & Thuật toán tính attention tối ưu memory --- giảm memory footprint đáng kể, cho phép dùng attention trên feature maps lớn. \\
\bottomrule
\end{longtable}

\section*{Kỹ thuật Detection}

\begin{longtable}{p{3.5cm}p{12cm}}
\toprule
\textbf{Thuật ngữ} & \textbf{Giải thích} \\
\midrule
\endfirsthead
\toprule
\textbf{Thuật ngữ} & \textbf{Giải thích} \\
\midrule
\endhead

\textbf{Bounding Box} & Hình chữ nhật bao quanh vật thể, xác định vị trí (x, y, width, height). \\
\midrule
\textbf{Anchor Box} & Hộp mẫu (template) được định nghĩa trước. Model dự đoán offsets so với anchor thay vì tọa độ tuyệt đối. v2--v7 dùng anchor, v8+ bỏ anchor. \\
\midrule
\textbf{Anchor-free} & Không dùng anchor. Dự đoán trực tiếp center point + 4 distances (khoảng cách từ tâm đến 4 cạnh). Đơn giản hơn anchor-based. \\
\midrule
\textbf{NMS} & Non-Maximum Suppression --- thuật toán loại bỏ bounding boxes trùng lặp. Giữ box confidence cao nhất, bỏ các boxes có IoU lớn. \\
\midrule
\textbf{NMS-free} & Không cần NMS --- model tự output 1 box/vật. Từ v10 trở đi. Dùng dual assignment (o2m + o2o). \\
\midrule
\textbf{Decoupled Head} & Tách classification và regression thành 2 branches riêng biệt. YOLOX giới thiệu, v6, v8+ adopt. \\
\midrule
\textbf{Label Assignment} & Chiến lược gán prediction cho ground truth khi training. SimOTA (động), TAL (task-aligned), STAL (small-target-aware). \\
\midrule
\textbf{End-to-End} & Pipeline từ input đến output không cần bước xử lý thủ công trung gian (ví dụ: không cần NMS). \\
\midrule
\textbf{Multi-scale} & Dự đoán ở nhiều kích thước (scales). YOLOv3+: 3 scales (lớn/trung/nhỏ) để detect vật ở mọi kích thước. \\
\midrule
\textbf{One-stage / Two-stage} & One-stage: 1 lần forward pass ra kết quả (YOLO). Two-stage: region proposal + classification (Faster R-CNN). \\
\bottomrule
\end{longtable}

\section*{Metrics (Chỉ số đánh giá)}

\begin{longtable}{p{3.5cm}p{12cm}}
\toprule
\textbf{Thuật ngữ} & \textbf{Giải thích} \\
\midrule
\endfirsthead
\toprule
\textbf{Thuật ngữ} & \textbf{Giải thích} \\
\midrule
\endhead

\textbf{IoU} & Intersection over Union --- đo mức trùng lặp giữa predicted box và ground truth. $IoU = \frac{Giao}{Hop}$. Từ 0 (không trùng) đến 1 (trùng hoàn hảo). \\
\midrule
\textbf{mAP} & mean Average Precision --- trung bình AP trên tất cả classes. Metric chính đánh giá detection. \\
\midrule
\textbf{mAP@50} & AP ở ngưỡng IoU $\geq$ 0.5. Dùng trên VOC. \\
\midrule
\textbf{mAP@50:95} & Trung bình AP ở IoU từ 0.5 đến 0.95 (bước 0.05). Metric chính trên COCO --- khắt khe hơn mAP@50. \\
\midrule
\textbf{FPS} & Frames Per Second --- số ảnh xử lý mỗi giây. $\geq$30 FPS = real-time. \\
\midrule
\textbf{Params} & Số lượng tham số của model. Ví dụ: 3.2M = 3.2 triệu tham số. Ít params = model nhẹ. \\
\midrule
\textbf{FLOPs / GFLOPs} & Floating Point Operations --- số phép tính cần thực hiện. Đo computational cost. 1 GFLOPs = 1 tỷ phép tính. \\
\midrule
\textbf{Latency} & Thời gian xử lý 1 ảnh (ms). Khác FPS vì bao gồm cả pre/post-processing. \\
\bottomrule
\end{longtable}

\section*{Training (Huấn luyện)}

\begin{longtable}{p{3.5cm}p{12cm}}
\toprule
\textbf{Thuật ngữ} & \textbf{Giải thích} \\
\midrule
\endfirsthead
\toprule
\textbf{Thuật ngữ} & \textbf{Giải thích} \\
\midrule
\endhead

\textbf{Batch Normalization} & Chuẩn hóa output mỗi layer theo mini-batch $\rightarrow$ training ổn định hơn, nhanh hội tụ hơn. YOLOv2 thêm BN: +2\% mAP. \\
\midrule
\textbf{Mosaic} & Data augmentation ghép 4 ảnh thành 1 $\rightarrow$ model thấy nhiều context, giảm batch size cần thiết. YOLOv4 giới thiệu. \\
\midrule
\textbf{CutMix} & Cắt 1 vùng từ ảnh A dán vào ảnh B, mix labels theo tỷ lệ diện tích. \\
\midrule
\textbf{Augmentation} & Biến đổi ảnh khi training (xoay, lật, zoom, đổi màu...) để model tổng quát hơn. \\
\midrule
\textbf{Transfer Learning} & Dùng model đã train trên dataset lớn (ImageNet) $\rightarrow$ fine-tune trên dataset nhỏ của mình. Tiết kiệm thời gian, ít data hơn. \\
\midrule
\textbf{Quantization} & Giảm precision từ FP32 (32-bit) xuống INT8 (8-bit) $\rightarrow$ model nhẹ hơn 4$\times$, nhanh hơn, nhưng có thể mất accuracy. \\
\midrule
\textbf{PTQ} & Post-Training Quantization --- quantize sau khi train xong. Nhanh nhưng mất accuracy. \\
\midrule
\textbf{QAT} & Quantization-Aware Training --- train lại với simulated quantization $\rightarrow$ giữ accuracy tốt hơn PTQ. \\
\midrule
\textbf{BoF / BoS} & Bag of Freebies (tăng training cost, không tăng inference) / Bag of Specials (tăng nhẹ inference cost). YOLOv4 hệ thống hóa. \\
\bottomrule
\end{longtable}

\section*{Loss Functions (Hàm mất mát)}

\begin{longtable}{p{3.5cm}p{12cm}}
\toprule
\textbf{Thuật ngữ} & \textbf{Giải thích} \\
\midrule
\endfirsthead
\toprule
\textbf{Thuật ngữ} & \textbf{Giải thích} \\
\midrule
\endhead

\textbf{MSE / SSE} & Mean/Sum Squared Error --- đo sai số bình phương. v1--v3 dùng cho box coordinates. \\
\midrule
\textbf{IoU Loss} & $\mathcal{L} = 1 - IoU$. Trực tiếp tối ưu IoU thay vì tọa độ riêng lẻ. \\
\midrule
\textbf{GIoU} & Generalized IoU --- xét cả vùng bao ngoài 2 boxes. Giải quyết trường hợp 2 box không giao nhau (IoU = 0 nhưng GIoU $\neq$ 0). \\
\midrule
\textbf{CIoU} & Complete IoU --- thêm khoảng cách tâm + aspect ratio vào GIoU. $\mathcal{L}_{CIoU} = 1 - IoU + \frac{\rho^2}{c^2} + \alpha v$. YOLOv4+ dùng. \\
\midrule
\textbf{DFL} & Distribution Focal Loss --- dự đoán phân phối xác suất cho box coordinates (16 bins) thay vì 1 giá trị. v8--v12 dùng, YOLO26 bỏ. \\
\midrule
\textbf{BCE} & Binary Cross-Entropy --- loss cho classification. v3+ dùng BCE thay Softmax để hỗ trợ multi-label. \\
\bottomrule
\end{longtable}

\section*{Deployment (Triển khai)}

\begin{longtable}{p{3.5cm}p{12cm}}
\toprule
\textbf{Thuật ngữ} & \textbf{Giải thích} \\
\midrule
\endfirsthead
\toprule
\textbf{Thuật ngữ} & \textbf{Giải thích} \\
\midrule
\endhead

\textbf{ONNX} & Open Neural Network Exchange --- format chuẩn để export model, chạy được trên nhiều frameworks/hardware. \\
\midrule
\textbf{TensorRT} & NVIDIA SDK tối ưu inference trên GPU NVIDIA. Tăng 2--5$\times$ speed so với PyTorch. \\
\midrule
\textbf{Edge deployment} & Triển khai model trên thiết bị nhỏ (mobile, camera, drone) thay vì server lớn. Cần model nhẹ, ít computation. \\
\midrule
\textbf{FP16 / FP32} & Float 16-bit / 32-bit precision. FP16 nhanh hơn, ít memory hơn, gần như không mất accuracy. \\
\midrule
\textbf{INT8} & Integer 8-bit. Nhẹ hơn FP32 gấp 4$\times$ nhưng có thể mất 0.5--2\% accuracy. \\
\bottomrule
\end{longtable}
